\section{Preliminaries}

In this section, we introduce the probabilistic framework, the functional setting, and some useful inequalities that will be used throughout the paper.

\subsection{Probabilistic Framework}

Let

\[
(\Omega_P,\mathcal F,\mathbb P)
\]

be a complete probability space endowed with a filtration

\[
\{\mathcal F_t\}_{t\ge0},
\]

satisfying the usual conditions of completeness and right continuity.

We consider a standard one-dimensional Wiener process

\[
W=\{W_t\}_{t\ge0}
\]

adapted to the filtration

\[
\{\mathcal F_t\}_{t\ge0}.
\]

Throughout this work, all stochastic processes are assumed to be adapted to this filtration.

\subsection{Functional Spaces}

Let $\Omega\subset\mathbb R^2$ be a bounded polygonal domain.

We introduce the Hilbert space

\[
H=L^2(\Omega),
\]

equipped with the inner product

\[
(u,v)
=
\int_\Omega uv\,dx
\]

and norm

\[
\|u\|
=
(u,u)^{1/2}.
\]

We also define

\[
V=H_0^1(\Omega),
\]

equipped with the norm

\[
\|v\|_V
=
\|\nabla v\|_{L^2(\Omega)}.
\]

The Gelfand triple

\[
V
\subset
H
\subset
V'
\]

plays a central role in the variational formulation of the stochastic problem.

\subsection{Bochner Spaces}

For a Banach space $X$, we denote by

\[
L^2(0,T;X)
\]

the space of square-integrable functions from $(0,T)$ into $X$.

Similarly,

\[
L^2(\Omega_P;X)
\]

denotes the space of square-integrable random variables taking values in $X$.

We will frequently use the space

\[
L^2
\bigl(
\Omega_P;
L^2(0,T;V)
\bigr).
\]

\subsection{Poincaré Inequality}

The following inequality holds for all

\[
v\in H_0^1(\Omega).
\]

:contentReference[oaicite:0]{index=0}

where $C_P>0$ denotes the Poincaré constant.

\subsection{Young's Inequality}

For any

\[
a,b\in\mathbb R
\]

and any

\[
\varepsilon>0,
\]

the following inequality holds.

:contentReference[oaicite:1]{index=1}

This inequality will be used repeatedly in the derivation of energy estimates.

\subsection{Grönwall's Lemma}

We recall the classical continuous Grönwall lemma.

\begin{lemma}
Let $y:[0,T]\rightarrow\mathbb R$ be a nonnegative function satisfying

\[
y(t)
\le
C
+
\int_0^t
\alpha(s)y(s)\,ds,
\]

where

\[
C\ge0
\]

and

\[
\alpha\in L^1(0,T).
\]

Then

\[
y(t)
\le
C
\exp
\left(
\int_0^t
\alpha(s)\,ds
\right).
\]

\end{lemma}

The discrete version of Grönwall's lemma will also be used in the convergence analysis of the fully discrete finite element scheme.